\section{Evaluation}\label{sec:evaluation}

The approach was evaluated on the verification tasks of SV-COMP and on two benchmarks from the literature, with
BenchExec, the limits of SV-COMP for CPU time, 7\,GB of memory and 2 cores. The results of the other tools are the
official ones of SV-COMP~2026. \cuvee{} writes witnesses, but empty ones for now: witnesses are not really supported,
so its results are not confirmed, and no scores are given. Table~\ref{tab:evaluation} counts the correct answers by
the expected verdicts: \texttt{TEQ} is the engine alone, \texttt{TEQ+Z3} lowers the goals that remain for Z3
(\texttt{-fallback}), and \texttt{CHC} encodes the same programs as Horn clauses for Z3, a baseline from the same
frontend.

\begin{table}[t]
\centering\small
\setlength{\tabcolsep}{4pt}
\begin{tabular}{@{}lrrrrrrrrrr@{}}
\toprule
 & \multicolumn{2}{c}{tasks} & \multicolumn{2}{c}{\texttt{TEQ}} & \multicolumn{2}{c}{\texttt{TEQ+Z3}} &
   \multicolumn{2}{c}{\texttt{CHC}} & \multicolumn{2}{c}{best other} \\
 & true & false & true & false & true & false & true & false & true & false \\
\midrule
SV-COMP Arrays & 320 & 120 & 138 & 9 & 148 & 48 & 8 & 45 & 68 & 79 \\
SV-COMP Loops & 555 & 209 & 149 & 16 & 230 & 20 & 150 & 41 & 405 & 172 \\
MonoCera tests & 125 & -- & 79 & -- & 88 & -- & 0 & -- & 41 & -- \\
SV-LIB CTranslated & 190 & 59 & 140 & 29 & 142 & 29 & 129 & 40 & 118 & 38 \\
Rapid benchmarks & 33 & -- & 17 & -- & 17 & -- & 0 & -- & 30 & -- \\
\bottomrule
\end{tabular}
\caption{Correct answers, true and false, by the expected verdicts of the tasks. None of the configurations gave a
wrong answer. The best other tool is Symbiotic on Arrays, AISE on Loops, MonoCera on its own tests, PySvLib-CHC on
SV-LIB, and Rapid on its own benchmarks.}\label{tab:evaluation}
\end{table}

\paragraph{Arrays of unbounded size are the strength.} On ReachSafety-Arrays, lookup proves 138 safe tasks alone and
148 with Z3, against 68 for the best participant of SV-COMP~2026 (VeriAbs, which takes part out of competition, proves
more). These tasks fill arrays with values that depend on the position or on earlier values, over symbolic sizes; they
need invariants quantified over the array indices, and the summaries of the database do not depend on the number of
iterations. On the aggregates of the MonoCera tests (Section~\ref{sec:monocera}), lookup proves 79 of 125 and 88 with
Z3, against 41 for MonoCera, which instruments the programs for these properties.

\paragraph{Numeric loops and refutations are the weakness.} On ReachSafety-Loops, the approach stays behind the tools
based on abstract interpretation and model checking: many of its tasks need non-linear invariants over integers
(\texttt{nla-digbench}), where Z3 adds most, and many have loops with constant bounds, which bounded model checking
unrolls completely while lookup gains nothing from a small bound. Refutations are found only when a goal without
traces evaluates to false, or by Z3 on the remaining goals; the engine builds no model. Lookup and Horn clauses are
complementary: each proves tasks that the other does not.

\paragraph{The lemmas are reused.} Each proof reports the lemmas it applies (\texttt{-rules}). In the 286 proofs of
\texttt{TEQ} on Arrays and Loops, 109 of the 153 lemmas are applied, ten of them in at least 100 tasks; the lemmas
about counters and the lookup of $\stores$ carry most proofs. Of the 44 unused lemmas, 23 belong to sorting and search,
written for the algorithms of this document, which are not part of the benchmarks.

\paragraph{Annotations per task, for comparison.} For the 790 safe tasks of Arrays and Loops, a coding agent wrote ACSL
invariants and contracts for each task, with which a verification condition generator proves 711; lookup proves 285 of
these without annotations. Annotating each task this way is more successful, but the effort does not carry over to
other programs, whereas the lemmas are written once, reused, and proved.
